\section{Pasabajos}

A partir de la alimentación del circuito con una señal cuadrada de frecuencia 28,2 kHz ($\frac{f_0}{10}$), se obtuvo la siguiente salida en el capacitor dada por la figura~\ref{fig:4_0.png}. 
\begin{figure}[h]
        \centering
        \includegraphics[width=0.5\textwidth]{4_0.png}
        \caption{Verde: Salida de filtro pasabajos ($V_C$) a 28,2 kHz; Amarillo: señal de entrada ($V_{in}$) desde el OpAmp. Escala en tiempo: 10 $\mu s$/div; Escala en tensión: 500 mV/div.}
        \label{fig:4_0.png}
    \end{figure}



% ********* ITEM A *********
\subsection{Respuesta al Escalón}
En base a la medición realizada en la figura \ref{fig:4_0.png}, se obtuvo la siguiente respuesta al escalón: frecuencia de oscilación del transitorio: 238 kHz; tiempo de establecimiento del 5$\%$: $\approx 3,24 \mu$s; sobrepico de 1.38 V. 
Luego, los valores teóricos con los que contrastar los obtenidos experimentalmente son: frecuencia de oscilación del transitorio: 244,5 kHz; tiempo de establecimiento del 5$\%$: $\approx 3,31 \mu$s; sobrepico de 1.31 V. Llegamos a una diferencia de frecuencias de oscilación del transitorio del 2.6\%. La diferencia se puede deber a la tolerancia del capacitor, ya que es el componente de mayor tolerancia (20\%).



% ********* ITEM B *********
\subsection{Salida con Barrido de Frecuencias}
Variando la frecuencia de la señal de entrada, se obtuvieron las siguientes salidas observables en la figura~\ref{fig:comparativa_frecuencias}.
\begin{figure}[h]
    \centering
    % Primera fila
    \begin{subfigure}[b]{0.48\textwidth}
        \centering
        \includegraphics[width=\textwidth]{4b_01.png}
        \caption{Frecuencia: 28,2 kHz}
        \label{fig:f28k}
    \end{subfigure}
    \hfill
    \begin{subfigure}[b]{0.48\textwidth}
        \centering
        \includegraphics[width=\textwidth]{4b_02.png}
        \caption{Frecuencia: 65 kHz}
        \label{fig:f65k}
    \end{subfigure}

    \vspace{0.5cm} % Espacio entre filas

    % Segunda fila
    \begin{subfigure}[b]{0.48\textwidth}
        \centering
        \includegraphics[width=\textwidth]{4b_04.png}
        \caption{Frecuencia: 130 kHz}
        \label{fig:f130k}
    \end{subfigure}
    \hfill
    \begin{subfigure}[b]{0.48\textwidth}
        \centering
        \includegraphics[width=\textwidth]{4b_05.png}
        \caption{Frecuencia: 200 kHz}
        \label{fig:f200k}
    \end{subfigure}

    \caption{Comparación de la respuesta temporal del filtro para distintas frecuencias de entrada. Se observa la atenuación progresiva de la amplitud a medida que aumenta la frecuencia (comportamiento pasabajos). Escala en tiempo: 5$\mu$ s(a), 3,9 $\mu$ s(b), 1,84 $\mu$ s(c) y (d); Escala en tensión: 500 mV/div.}
    \label{fig:comparativa_frecuencias}
\end{figure}
Vemos que a medida que aumenta la frecuencia de la señal de entrada, la amplitud de la salida disminuye, lo cual es consistente con el comportamiento esperado de un filtro pasabajos. Viendo la situación del punto de vista gráfico, se observa que a medida que aumenta la frecuencia, el tiempo de establecimiento disminuye, hasta tener una salida totalmente oscilante. 



% ********* ITEM C *********
\subsection{Cálculo de la Función Transferencia}

A partir del circuito de la figura~\ref{fig:circuito_pasabajos}, la función 
transferencia se obtiene mediante un divisor de impedancias, tomando como 
salida la tensión en el capacitor:

\begin{equation}
    H(s) = \frac{V_o(s)}{V_i(s)} = \frac{Z_C}{Z_L + Z_R + Z_C} 
         = \frac{\dfrac{1}{sC}}{sL + R + \dfrac{1}{sC}}
    \label{eq:H_divisor}
\end{equation}

Multiplicando numerador y denominador por $sC$:

\begin{equation}
    H(s) = \frac{1}{s^2LC + sRC + 1}
    \label{eq:H_s}
\end{equation}

Reescribiendo en forma canónica de segundo orden, con frecuencia natural 
$\omega_0 = \frac{1}{\sqrt{LC}}$ y coeficiente de amortiguamiento 
$\xi = \frac{R}{2}\sqrt{\frac{C}{L}}$:

\begin{equation}
    H(s) = \frac{\omega_0^2}{s^2 + 2\xi\omega_0\, s + \omega_0^2}
    \label{eq:H_canonica}
\end{equation}

Con los valores de los componentes: ($L = \SI{38.3}{\micro\henry}$, 
$C = \SI{8.58}{\nano\farad}$, $R = \SI{68}{\ohm}$), los parámetros resultan:

\begin{equation}
    f_0 = \frac{1}{2\pi\sqrt{LC}} 
        = \frac{1}{2\pi\sqrt{38{,}3\times10^{-6}\cdot 8{,}58\times10^{-9}}} 
        \approx \SI{277.8}{\kilo\hertz}
    \label{eq:f0_pasabajos}
\end{equation}

\begin{equation}
    \xi = \frac{R}{2}\sqrt{\frac{C}{L}} 
        = \frac{68}{2}\sqrt{\frac{8{,}58\times10^{-9}}{38{,}3\times10^{-6}}} 
        \approx 0{,}508
    \label{eq:xi_pasabajos}
\end{equation}

Dado que $\xi < \frac{1}{\sqrt{2}}$, el sistema es subamortiguado y presenta 
sobrepico tanto en la respuesta al escalón como en la respuesta en frecuencia.



% ********* ITEM D *********
\subsection{Diagrama de Bode Teórico-Práctico}
\begin{figure}[H]
        \centering
        \includegraphics[width=0.7\textwidth]{4d_opam.png}
        \caption{Verde: Salida de filtro pasabajos ($V_C$) a 28,2 kHz; Amarillo: señal de entrada ($V_{in}$) desde el OpAmp. Escala en tiempo: 10 $\mu s$/div; Escala en tensión: 500 mV/div.}
        \label{fig:4d_opam.png}
    \end{figure}


\subsubsection{R = 0 \texorpdfstring{$\Omega$}{pi} }



%%%%%%%%%%%%%%%%%%%%%%%%%%%%%%%%%%%%%%%%%%%
%%%%%%%%%%%%%%%%%%%%%%%%%%%%%%%%%%%%%%%%%%%
%%%%%%%%%%%%%%%%%%%%%%%%%%%%%%%%%%%%%%%%%%%
%   ESTO NO SE SI HACE FALTA              %
%%%%%%%%%%%%%%%%%%%%%%%%%%%%%%%%%%%%%%%%%%%
%%%%%%%%%%%%%%%%%%%%%%%%%%%%%%%%%%%%%%%%%%%
%%%%%%%%%%%%%%%%%%%%%%%%%%%%%%%%%%%%%%%%%%%
\subsection{Respuesta a la Señal Cuadrada}

se uso una señal cuadrada de $f = f_0/10 \approx \SI{27.8}{\kilo\hertz}$ y amplitud 
de $v_{o,\,\text{max}} = \SI{2}{V_{pp}}$. como el periodo de la entrada ($T = \SI{36}{\micro\second}$) es 
diez veces el de oscilacion propia ($T_0 = \SI{3.6}{\micro\second}$), el transitorio llega a terminar en 
cada semiciclo antes de que venga el proximo flanco. por esto, en la salida $v_o(t)$ se ve bien la respuesta
 al escalon y el comportamiento subamortiguado en cada transicion.


Ante un escalón de amplitud $A$, la respuesta temporal de un sistema de segundo 
orden subamortiguado es:

\begin{equation}
    v_o(t) = A\left[1 - \frac{e^{-\xi\omega_0 t}}{\sqrt{1-\xi^2}}
             \sin\!\left(\omega_d t + \phi\right)\right], 
             \quad \phi = \arccos(\xi)
    \label{eq:respuesta_escalon}
\end{equation}

donde la frecuencia de oscilación del transitorio es:

\begin{equation}
    \omega_d = \omega_0\sqrt{1-\xi^2} \implies 
    f_d = f_0\sqrt{1-\xi^2} 
        = \SI{277.8}{\kilo\hertz}\cdot\sqrt{1-0{,}508^2} 
        \approx \SI{240.8}{\kilo\hertz}
    \label{eq:fd}
\end{equation}

El sobrepico máximo, que ocurre en $t_p = \pi/\omega_d$, vale:

\begin{equation}
    M_P = e^{-\pi\xi/\sqrt{1-\xi^2}} 
        = e^{-\pi \cdot 0{,}508/\sqrt{1-0{,}508^2}} 
        \approx 0{,}163
    \label{eq:MP}
\end{equation}

Es decir, la salida supera el valor final en un $16{,}3\%$ del escalón aplicado.